% Source: projections_adjoin_and_ls_geometry.tex
\begin{notechapter}
\chapter{Projections, Adjoints and the Geometry of Least Squares}\label{ch:least-squares-geometry}

\section*{Introduction}
This document synthesizes key concepts in linear algebra regarding projections, the role of the adjoint operator, and their implications for solving linear systems and understanding Machine Learning techniques like Least Squares (LS) and Principal Component Analysis (PCA). We clarify the distinction between a general mapping matrix $A$ and a specific projection matrix $P$.

\section{Projection Matrices: Idempotence and Symmetry}

A linear transformation $P: V \to V$ is a \textbf{projection operator} if it is \textbf{idempotent}, meaning applying the projection twice is the same as applying it once:
$$P^2 = P$$
The range of $P$, denoted $\text{Im}(P)$, is the subspace onto which we are projecting. \textbf{Not all matrices are projection matrices; they must satisfy this $P^2=P$ condition.}

\subsection{The Role of the Adjoint}

The adjoint operator $P^*$ (conjugate transpose, or simply the transpose $P^T$ in real vector spaces) is defined by the inner product relation:
$$\langle Pv, w \rangle = \langle v, P^*w \rangle$$
This operator is crucial for distinguishing between types of projections:

\begin{enumerate}[label=(\roman*)]
    \item \textbf{Orthogonal Projection:} A projection $P$ is \textbf{orthogonal} if it is \textbf{self-adjoint}, $P^* = P$. Orthogonal projection matrices are thus both idempotent and symmetric (or Hermitian in the complex case).
    \item \textbf{Oblique Projection:} If $P^* \neq P$, the projection is oblique. The error vector is not orthogonal to $\text{Im}(P)$.
\end{enumerate}

\section{Least Squares and the Adjoint of a General Matrix $A$}

In Machine Learning and statistics, we work with a general mapping matrix $A$ (which is typically \textbf{not} a projection matrix itself) to derive optimal orthogonal projections. Consider a data matrix $A$ mapping from an input space $\mathbb{R}^n$ to an output space $\mathbb{R}^m$.

\subsection{Mapping Between Spaces}
The operators $A$ and its transpose $A^T$ map between the input and output spaces:
\begin{align*}
A: \mathbb{R}^n &\to \mathbb{R}^m \\
A^T: \mathbb{R}^m &\to \mathbb{R}^n
\end{align*}
The transpose $A^T$ takes vectors from the output space and maps them back into the input space. This $A^T$ is the matrix representation of the adjoint operator, relying on the standard dot product (inner product) in $\mathbb{R}^n$ and $\mathbb{R}^m$.

\subsection{Adjoint vs. Dual Map (Clarification)}
In practical applications like Least Squares, we are concerned with geometry and orthogonality, which are inner product concepts. Therefore, $A^T$ always refers to the matrix of the \textbf{adjoint operator} in this context. The standard inner product allows the matrix of the adjoint to be simply the algebraic transpose.

\subsection{The Normal Equations}

When solving the system $A\hat{x} = y$ in a least squares sense (because $y$ is not in the column space of $A$, $\text{Im}(A)$), we seek the $\hat{x}$ that minimizes the orthogonal error $e = y - A\hat{x}$. The adjoint condition ensures this error is orthogonal to the image space\footnote{This relationship is formally proven in Appendix A, showing that $\text{Ker}(A^T)$ is the orthogonal complement of $\text{Im}(A)$.}:
$$A^T e = 0 \implies A^T(y - A\hat{x}) = 0 \implies A^TA\hat{x} = A^Ty$$
These are the \textbf{normal equations}. We solve this system for the coefficients $\hat{x}$, which define the best linear combination of the columns of $A$. The resulting vector $A\hat{x}$ is the orthogonal projection of $y$ onto $\text{Im}(A)$.

The orthogonal projection of $y$ onto $\text{Im}(A)$ always exists and is unique, even when $A^TA$ is not invertible. If $A^TA$ is invertible, the projection matrix is $P_{\text{proj}} = A(A^TA)^{-1}A^T$. If $A^TA$ is not invertible, we use the Moore-Penrose pseudoinverse to calculate the projection matrix and the minimum-norm solution for $\hat{x}$ (see Appendix B).

\section{Existence and Uniqueness of Solutions}

The existence and uniqueness of the solution $\hat{x}$ to the system $A^TA\hat{x} = A^Ty$ depends entirely on the properties of the kernel of $A$ ($\text{Ker}(A)$) and the image of $A$ ($\text{Im}(A)$).

\begin{itemize}[leftmargin=*]
    \item \textbf{Unique Solution:} A unique $\hat{x}$ exists if and only if the columns of $A$ are linearly independent ($\text{Ker}(A) = \{0\}$). In this case, $A^TA$ is invertible, and we have $\hat{x} = (A^TA)^{-1}A^Ty$. This occurs when $A$ has more rows than columns ($m \geq n$) and is of full rank.
    \item \textbf{Infinite Solutions:} If the columns of $A$ are linearly dependent ($\text{Ker}(A) \neq \{0\}$), $A^TA$ is not invertible, and there are infinitely many $\hat{x}$ that achieve the same minimum error. This is guaranteed when $A$ has more columns than rows ($n > m$).
\end{itemize}

\subsection*{Summary of Cases for $A\hat{x}=y$}
The system $A\hat{x}=y$ and its corresponding least squares solution fall into four scenarios:

\begin{table}[h!]
    \centering
    \begin{tabular}{|l|l|l|}
    \hline
    & $y \in \text{Im}(A)$ (Exact Solution Exists) & $y \notin \text{Im}(A)$ (LS Needed) \\
    \hline
    $\text{Ker}(A) = \{0\}$ (Unique Fit) & Unique solution for $x$. & Unique LS solution $\hat{x}$. \\
    \hline
    $\text{Ker}(A) \neq \{0\}$ (Redundancy) & Infinite solutions for $x$. & Infinite LS solutions for $\hat{x}$. \\
    \hline
    \end{tabular}
\end{table}

\section{Geometric Visualization of Least Squares}

We can visualize the least squares problem in $\mathbb{R}^3$, projecting onto a 2D plane (the image of $A$).

\begin{figure}[h!]
    \centering
    \begin{tikzpicture}[scale=2.0, 
        % Define a view angle (rotate around Z, then X) for better perspective
        x={(-0.5cm,-0.3cm)}, 
        y={(1cm,-0.2cm)}, 
        z={(0cm,1cm)}
    ]
            % Define coordinates based on the new perspective axes
            \coordinate (O) at (0,0,0);
            \coordinate (X_axis) at (3,0,0);
            \coordinate (Y_axis) at (0,3,0);
            \coordinate (Z_axis) at (0,0,3);
            
            % The plane (Im(A)) - a non-axial plane for better generalization
            % Define corners of the diamond plane shape
            \coordinate (P1) at (2.5,0,0);
            \coordinate (P2) at (0,2.5,0);
            \coordinate (P3) at (-1.0,0,0);
            \coordinate (P4) at (0,-1.0,0);
            
            \fill[blue!15, opacity=0.7] (P1) -- (P2) -- (P3) -- (P4) -- cycle;
            \draw[blue, dashed] (P1) -- (P2) -- (P3) -- (P4) -- cycle;
            \draw[blue, thick] (P1) -- (P2);
            \node[blue] at (1.0, 1.0, -0.5) {$\text{Im}(A)$};

            % The vector y (set coordinates to be above the center of the plane area)
            \coordinate (y) at (0.5, 0.7, 2.5);
            \draw[->, very thick, black] (O) -- (y) node[above right] {$y$};

            % The projection of y (A*x_hat) - point in the plane
            \coordinate (Ahatx) at (0.5, 0.7, 0);
            \draw[->, thick, red] (O) -- (Ahatx) node[below left] {$A\hat{x}$};

            % The error vector (e) - orthogonal to the plane
            \draw[->, thick, dashed, green!60!black] (Ahatx) -- (y) node[right, midway] {$e$};
            % Add label specifically for 'e' to avoid overlap
            %\node[green!60!black, right, font=\small] at ($(Ahatx)!0.5!(y)$) {Error};

            % Right angle marker
            % Calculate points for a clear right angle marker
            \draw ($(Ahatx)!0.1cm!(y)$) -- ($($(Ahatx)!0.15cm!(y)$)!0.09cm!90:(y)$) -- ($(Ahatx)!0.12cm!(O)$);

            % Axes (drawing them behind the main action if possible for clarity)
            \draw[->, gray!50, dashed] (O) -- (X_axis) node[below left] {$x_1$};
            \draw[->, gray!50, dashed] (O) -- (Y_axis) node[below right] {$x_2$};
            \draw[->, gray!50, dashed] (O) -- (Z_axis) node[above left] {$x_3$};
        \end{tikzpicture}
    \caption{Geometric visualization of Least Squares. The vector $y$ (in $\mathbb{R}^3$) is projected orthogonally onto the subspace $\text{Im}(A)$ (the plane) to find the closest point $A\hat{x}$. The error vector $e$ is perpendicular to the plane.}
\end{figure}

\chapterappendix
\section{Proof of the Fundamental Orthogonality Relation}

Here we provide the formal proof using the adjoint definition for the fundamental theorem of linear algebra relationship that the kernel of $A^T$ is the orthogonal complement of the image of $A$, i.e., $\text{Ker}(A^{T})=(\text{Im}(A))^{\perp }$.

Let $A: \mathbb{R}^n \to \mathbb{R}^m$ be a linear mapping represented by an $m \times n$ matrix $A$. The adjoint $A^T$ maps $\mathbb{R}^m \to \mathbb{R}^n$ and satisfies the inner product relationship $\langle Av, w \rangle = \langle v, A^Tw \rangle$ for all $v \in \mathbb{R}^n, w \in \mathbb{R}^m$.

We aim to prove the set equality $\text{Ker}(A^T) = (\text{Im}(A))^\perp$ by showing two inclusions:

\subsubsection*{Part 1: $\text{Ker}(A^T) \subseteq (\text{Im}(A))^\perp$}

Let $z \in \text{Ker}(A^T)$. By definition, $A^T z = 0$.
We want to show that $z$ is orthogonal to every vector in $\text{Im}(A)$. Any vector in $\text{Im}(A)$ can be written as $y = Ax$ for some $x \in \mathbb{R}^n$.

Consider the inner product of $y$ and $z$:
$$ \langle y, z \rangle = \langle Ax, z \rangle $$
Using the definition of the adjoint operator, we can move the $A$ to the other side as $A^T$:
$$ \langle Ax, z \rangle = \langle x, A^Tz \rangle $$
Since we know $z \in \text{Ker}(A^T)$, we substitute $A^Tz = 0$:
$$ \langle x, 0 \rangle = 0 $$
Thus, $\langle y, z \rangle = 0$ for all $y \in \text{Im}(A)$. This proves that $z$ is in the orthogonal complement of $\text{Im}(A)$, so $\text{Ker}(A^T) \subseteq (\text{Im}(A))^\perp$.

\subsubsection*{Part 2: $(\text{Im}(A))^\perp \subseteq \text{Ker}(A^T)$}

Let $w \in (\text{Im}(A))^\perp$. By definition, $w$ is orthogonal to every vector in $\text{Im}(A)$. This means for all $v \in \mathbb{R}^n$, $\langle Av, w \rangle = 0$.

Again using the definition of the adjoint:
$$ \langle Av, w \rangle = \langle v, A^Tw \rangle $$
We have $\langle v, A^Tw \rangle = 0$ for all possible vectors $v \in \mathbb{R}^n$.
The only vector $u$ such that $\langle v, u \rangle = 0$ for *all* $v$ in the vector space is the zero vector itself ($u=0$).
Therefore, we must have $A^Tw = 0$.

By definition, $w$ is in the kernel of $A^T$, so $(\text{Im}(A))^\perp \subseteq \text{Ker}(A^T)$.

\subsubsection*{Conclusion}

Since we have proven both inclusions, the sets are equal:
$$ \text{Ker}(A^T) = (\text{Im}(A))^\perp $$

\section{The Moore-Penrose Pseudoinverse and SVD}

When $A^TA$ is not invertible (due to linearly dependent columns in $A$), the standard formula for the projection matrix $P_{\text{proj}} = A(A^TA)^{-1}A^T$ fails. We still have a unique orthogonal projection onto $\text{Im}(A)$, which can be found using the \textbf{Moore-Penrose pseudoinverse}, denoted $A^+$.

\subsubsection*{Computing $A^+$ using SVD}
The SVD factors any $m \times n$ matrix $A$ into three matrices:
$$A = U \Sigma V^T$$
\begin{itemize}
    \item $U$: An $m \times m$ orthogonal matrix of left singular vectors.
    \item $\Sigma$: An $m \times n$ rectangular diagonal matrix containing the singular values $\sigma_i$.
    \item $V$: An $n \times n$ orthogonal matrix of right singular vectors.
\end{itemize}

The pseudoinverse $A^+$ is calculated using the pseudoinverse of the diagonal matrix $\Sigma^+$, via the formula:
$$A^+ = V \Sigma^+ U^T$$
To form $\Sigma^+$, we transpose $\Sigma$ and replace every non-zero singular value $\sigma_i$ with its reciprocal $1/\sigma_i$. Zero values in $\Sigma$ remain zero in $\Sigma^+$.

\subsubsection*{Mathematical Link between U and V (Why it works)}

The structure $A = U \Sigma V^T$ inherently links the columns of $U$ and $V$. Let $u_i$ be the $i$-th column of $U$ and $v_i$ be the $i$-th column of $V$. The SVD equation means that for any non-zero singular value $\sigma_i$, we have the direct relationship:
$$A v_i = \sigma_i u_i$$
This shows that $A$ maps the orthonormal basis vector $v_i$ in the input space ($\mathbb{R}^n$) directly to the orthonormal basis vector $u_i$ in the output space ($\mathbb{R}^m$), scaled by $\sigma_i$.

Now we analyze the action of the pseudoinverse $A^+$ on the vectors $u_i$:
$$A^+ u_i = (V \Sigma^+ U^T) u_i$$
\begin{itemize}
    \item The term $U^T u_i$ results in the standard basis vector $e_i$ because $U$ is orthogonal.
    \item $\Sigma^+ e_i$ results in $e_i$ scaled by $1/\sigma_i$ (for non-zero $\sigma_i$).
    \item $V (e_i / \sigma_i)$ results in the vector $v_i / \sigma_i$.
\end{itemize}

The action of the pseudoinverse on the basis vectors of the image space is:
$$A^+ u_i = \frac{1}{\sigma_i} v_i$$
The mathematical relationship shows exactly how $V$ relates to $U$ in the reverse operation: $A$ sends $v_i$ to $\sigma_i u_i$, and $A^+$ sends $u_i$ to $(1/\sigma_i) v_i$. $A^+$ acts as a true inverse on the subspaces spanned by the $u_i$ and $v_i$ corresponding to non-zero singular values. This complementary relationship ensures the construction is mathematically sound and consistent with the definition of an inverse on the relevant subspaces, always producing the unique orthogonal projection matrix $P_{\text{proj}} = AA^+$.
\end{notechapter}
